%% ============================================================
%%  SECCIÓN 2 — Resumen
%%  Responsable: Vera Vivanco, Jesús Rodolfo
%% ============================================================

\section{Resumen}

El presente informe de laboratorio tuvo como objetivo la determinación del coeficiente de dilatación lineal ($\alpha$) para tres materiales de uso común en la industria y la construcción: aluminio, cobre y vidrio comercial. El experimento se llevó a cabo utilizando un dilatómetro térmico clásico, en el cual se indujo la expansión de tubos huecos de cada material mediante el flujo de vapor de agua. La temperatura del sistema se incrementó desde una temperatura ambiental inicial de 26$^\circ$C hasta una temperatura final de ebullición de 100$^\circ$C, lo que representó una variación térmica controlada de $\Delta T = 74\ ^\circ\text{C}$.

Para medir con precisión la variación de longitud submilimétrica ($\Delta L$), se empleó el método mecánico de la aguja giratoria. Este procedimiento aprovecha el rodamiento sin deslizamiento del extremo libre del tubo sobre una aguja cilíndrica de diámetro $D = \qty{0,60}{\milli\meter}$, amplificando el desplazamiento lineal en una rotación angular $\theta$ que es registrada por un indicador circular graduado. La relación geométrica utilizada para la determinación del alargamiento fue $\Delta L = 2R\Delta\theta$ (donde $R = \qty{0,30}{\milli\meter}$ es el radio de la aguja), a partir de la cual se estimaron los coeficientes experimentales y se propagaron sus incertidumbres asociadas mediante el método de suma en cuadratura.

Los coeficientes de dilatación lineal obtenidos experimentalmente para cada material con su respectiva incertidumbre y error relativo porcentual respecto a la literatura científica fueron:
\begin{itemize}
    \item Aluminio: $\alpha_{\text{Al}} = (2{,}46 \pm 0{,}11) \times 10^{-5}\ ^\circ\text{C}^{-1}$ con un error relativo de \qty{6.8}{\percent}.
    \item Cobre: $\alpha_{\text{Cu}} = (1{,}60 \pm 0{,}07) \times 10^{-5}\ ^\circ\text{C}^{-1}$ con un error relativo de \qty{5.8}{\percent}.
    \item Vidrio borosilicato: $\alpha_{\text{vidrio}} = (2{,}43 \pm 0{,}14) \times 10^{-6}\ ^\circ\text{C}^{-1}$ con un error relativo de \qty{1.25}{\percent}.
\end{itemize}
Los resultados obtenidos muestran una excelente concordancia con los valores teóricos de referencia, lo que convalida la metodología experimental implementada y la precisión del factor de amplificación geométrica por rodadura doble.